\begin{proposition}[Conductor and exceptional pole divisor]\hypertarget{tag-02YI}{}\label{night:P4R.6}\label{stmt:P4R.6}
\StatementRecord{night:P4R.6}{proved}{02YI}{\SeriesRef{night:P4.12}}

Let $f:Y\to B$ be a proper surjective holomorphic map with connected
fibres from a smooth compact complex threefold with $b_1(Y)=b_2(Y)=0$
to a smooth compact curve, and suppose that its smooth general fibre
is a complex two-torus. Let $mS$ be a fibre with irreducible, reduced,
singular support $S$, let $\nu:N\to S$ be its normalization, and let
$\pi:T\to N$ be its minimal resolution. Put $\rho=\nu\circ\pi$.
There is an effective integral divisor $D$ on $T$ such that
\[
 \omega_T(D)\simeq\rho^*\omega_S,
 \qquad \operatorname{Supp}D
       =\rho^{-1}(\operatorname{Sing}S)_{\rm red}.
\]
In particular $K_T+D$ is torsion, by \cref{night:P4.12}.
The strict-transform conductor has its usual positive integral
coefficients. Every connected exceptional configuration over a
singular point of $N$ has strictly positive coefficients in $D$.
Every such configuration meets the strict-transform conductor.
\end{proposition}
\begin{proof}
\begin{enumerate}[label=\arabic*.,leftmargin=*]
\item The reduced support is a Cartier hypersurface in the smooth
threefold, so $\omega_S$ is invertible. A local generator $\eta$
pulls back from $S_{\rm reg}$ to a meromorphic two-form on $T$.
Its divisorial orders are integers. Changing $\eta$ by a holomorphic
unit changes its pullback by a unit, so the negatives of these
orders define an integral divisor $D$ and the displayed
line-bundle isomorphism.

Write $\mathfrak c$ for the conductor, viewed as an ideal on $N$.
For a finite birational normalization $A\subset B$, let $Q(A)$
be the total quotient ring obtained by inverting the non-zero-divisors
of $A$. Then $B\otimes_A Q(A)=Q(A)$. Every $A$-linear map
$\varphi:B\to A$ extends to a $Q(A)$-linear endomorphism of this
common total quotient ring, hence to multiplication by $\varphi(1)$.
Thus evaluation at $1$ identifies $\operatorname{Hom}_A(B,A)$ with
$\{a\in A:aB\subset A\}$. Finite duality therefore gives
\[
 \nu_*\omega_N
   =\mathcal Hom_{\mathcal O_S}(\nu_*\mathcal O_N,\omega_S),
 \qquad
 \omega_N\simeq\mathfrak c\otimes\nu^*\omega_S.
\]
At a regular generic point of a conductor curve, write
$\mathfrak c=(z^c)$, with $c>0$. The pullback of $\eta$ is
$z^{-c}$ times a regular canonical generator. Thus the
strict-transform coefficients of $D$ are exactly the positive
integral conductor exponents. Outside the conductor and the
exceptional curves, $\rho$ identifies the forms on the smooth
loci, and the order is zero.

\item Let $E_1,\ldots,E_r$ be the exceptional curves over a singular
point $p\in N$, and put $M_{ij}=E_i\cdot E_j$.
The fibre $\pi^{-1}(p)$ is connected: a decomposition into two
disjoint compact parts would, by properness, give disjoint
neighborhoods containing $\pi^{-1}(U)$ for a sufficiently small
neighborhood $U$ of $p$. The function with values $0$ and $1$
on those neighborhoods descends to a bounded holomorphic function
on $U\setminus\{p\}$. Normality extends it across $p$, whereas
approach to the two exceptional parts gives different limiting
values. This is a contradiction.

Choose a nonzero local holomorphic function on $N$ vanishing at
$p$. Its pullback divisor is $H+\sum_i n_iE_i$, with $n_i>0$
and $H$ effective without exceptional components. The strict
transform $H$ meets the exceptional fibre by properness. Since
the divisor is principal near that fibre,
\[
 Mn=-h,\qquad h_i=H\cdot E_i\ge0,\qquad h\ne0.
\]
The off-diagonal entries of $M$ are nonnegative. With
$v_i=n_iw_i$, direct expansion gives
\[
 v^tMv
 =-\sum_{i<j}M_{ij}n_in_j(w_i-w_j)^2
   -\sum_i n_ih_iw_i^2.
\]
This is the weighted intersection identity of
\cite[p.~6]{Mumford1961}: the local principal divisor above
replaces the source's hyperplane difference, and
$v=\operatorname{diag}(n_i)w$ is the change of variables.
Connectedness and $h\ne0$ make the right-hand side strictly
negative for $v\ne0$.

For $Md=-b$, $b_i\ge0$, set $u_i=d_i/n_i$. The same divisor
equation gives the row identity
\[
 b_i=h_i u_i+\sum_{j\ne i}M_{ij}n_j(u_i-u_j).
\]
A negative minimum of the $u_i$ forces every neighboring value
to equal it and forces $h_i=0$. Propagation through the connected
graph reaches a vertex with $h_i>0$, a contradiction. Thus
$d_i\ge0$. When $b\ne0$, a zero $d_i$ forces all neighboring
$d_j$ to be zero, and propagation contradicts $Md=-b$.
Consequently $d_i>0$ in that case; equivalently, $-M^{-1}$ has
strictly positive entries on the connected block.

\item Write $D=D_{\rm str}+\sum_i d_iE_i$, where
$D_{\rm str}$ is the strict-transform conductor. Adjunction and
minimality of $\pi$ give
\[
 K_T\cdot E_i=2p_a(E_i)-2-E_i^2\ge0.
\]
Indeed $E_i^2<0$ by the preceding identity, and an integral curve
with $p_a(E_i)=0$ is smooth and rational. Minimality excludes an
exceptional curve of this kind with self-intersection $-1$.
Also $D_{\rm str}\cdot E_i\ge0$. Since $\rho(E_i)$ is a point,
the restriction of $\rho^*\omega_S$ to $E_i$ has degree zero.
The canonical divisor identity therefore gives
\[
 Md=-b,\qquad b_i=(K_T+D_{\rm str})\cdot E_i\ge0.
\]
It follows that $d_i\ge0$. A zero coefficient forces $b=0$
on the connected block. Each of its curves is then a smooth
rational $(-2)$-curve, and the block is disjoint from
$D_{\rm str}$.

\item We now prove that every exceptional block meets
$D_{\rm str}$. The identity in the first item makes
$\mathfrak c$ a rank-one reflexive ideal on the normal surface
$N$. Such an ideal is the intersection of its codimension-one
localizations in the fraction field. Hence
$\operatorname{Supp}(\mathcal O_N/\mathfrak c)$ has no isolated
point: a conductor ideal equal to $\mathcal O_N$ at all nearby
codimension-one points is equal to $\mathcal O_N$ itself.
A conductor curve through $p$ has strict transform meeting
$\pi^{-1}(p)$, by properness. Therefore a block disjoint from
$D_{\rm str}$ lies over a point where $\nu$ is an isomorphism.
Its image in $S$ would be a normal isolated hypersurface
singularity.

We exclude such a point directly. In a coordinate ball write
$f=u g^m$, where $g=0$ is the reduced equation and $u$ is a
holomorphic unit. Choose a logarithm of $u$ and put
$h=\exp((\log u)/m)g$. Then $f=h^m$ exactly. The critical point
of $h$ is isolated: a positive-dimensional critical component
through the point would have constant value zero and would lie
in the singular locus of the normal surface $\{h=0\}$.
The gradient of $h$ has positive local degree. A sufficiently
small regular value $a$ of that gradient makes
$h_a=h-\sum_j a_jz_j$ have finitely many, and at least one,
nondegenerate critical points. The degree and positivity
argument is given in \cite[Appendix~B, pp.~111--114]{Milnor1968}.
For this analytic germ the gradient is finite over a small
target neighborhood; at each regular preimage its real
Jacobian is the positive number $|\det_{\mathbb C}d(\nabla h)|^2$.

The fibre identifications can be made inside a fixed ball.
Indeed the critical locus of squared radius on
$\{h=0\}\setminus\{0\}$ is semianalytic. In a compatible locally
finite analytic triangulation from
\cite[Theorem~2, p.~464]{Lojasiewicz1964}, an accumulation at $0$
would give a connected positive-dimensional simplex with $0$
in its closure. Squared radius would be constant on it and
equal to zero, a contradiction. A small ball therefore has
boundary transverse to $\{h=0\}$. Transversality on this compact
boundary persists for small fibre values and small perturbations.
Choose a nonzero value $\delta$ and then $a$ small enough that
all critical values of $h-sa\cdot z$, $0\le s\le1$, have
absolute value less than $|\delta|/2$. Boundary-tangent
horizontal vector fields identify the fibres over $\delta$
through this proper family.

At a critical point of $h_a$, holomorphic Morse coordinates give
$h_a=c+z_1^2+z_2^2+z_3^2$. In the fibre over $c+\epsilon$,
the real sphere $\sum x_j^2=\epsilon>0$ has normal bundle
$J(TS^2)$. The tangent-then-normal order
$(e_1,e_2,Je_1,Je_2)$ has the opposite orientation to the
complex order $(e_1,Je_1,e_2,Je_2)$, so its self-intersection
is $-\chi(S^2)=-2$. Transport along a path avoiding the finite
critical values, followed by the preceding parameter-family
identification, embeds this sphere in $h^{-1}(\delta)$.
Since $f=h^m$, this local fibre embeds in the smooth global
fibre $f^{-1}(\delta^m)$, a complex two-torus. The sphere lifts
to its contractible universal cover $\mathbb C^2$, so its
homology class and self-intersection in the torus are zero.
Its normal bundle is unchanged by the inclusion of the local
fibre, contradicting the value $-2$.

Thus every exceptional block meets $D_{\rm str}$, so its vector
$b$ in the third item is nonzero. All its coefficients $d_i$
are strictly positive. Every normal point of $S$ is smooth by
the same local argument. The singular locus is therefore the
conductor locus, and its inverse image under $\rho$ consists
exactly of the strict conductor and these exceptional fibres.
This proves effectiveness and the support equality.
Finally, \cref{night:P4.12} gives $\omega_S^{\otimes m}\simeq
\mathcal O_S$, whence
\[
 (\omega_T(D))^{\otimes m}
   \simeq\rho^*(\omega_S^{\otimes m})\simeq\mathcal O_T.
\]
\end{enumerate}
\end{proof}